\IfFileExists{stacks-project.cls}{%
\documentclass{stacks-project}
}{%
\documentclass{amsart}
}

% For dealing with references we use the comment environment
\usepackage{verbatim}
\newenvironment{reference}{\comment}{\endcomment}
%\newenvironment{reference}{}{}
\newenvironment{slogan}{\comment}{\endcomment}
\newenvironment{history}{\comment}{\endcomment}

% For commutative diagrams we use Xy-pic
\usepackage[all]{xy}

% We use 2cell for 2-commutative diagrams.
\xyoption{2cell}
\UseAllTwocells

% We use multicol for the list of chapters between chapters
\usepackage{multicol}

% This is generally recommended for better output
\usepackage{lmodern}
\usepackage[T1]{fontenc}

% For cross-file-references
\usepackage{xr-hyper}

% Package for hypertext links:
\usepackage{hyperref}

% For any local file, say "hello.tex" you want to link to please
% use \externaldocument[hello-]{hello}
\externaldocument[introduction-]{introduction}
\externaldocument[conventions-]{conventions}
\externaldocument[sets-]{sets}
\externaldocument[categories-]{categories}
\externaldocument[topology-]{topology}
\externaldocument[sheaves-]{sheaves}
\externaldocument[sites-]{sites}
\externaldocument[stacks-]{stacks}
\externaldocument[fields-]{fields}
\externaldocument[algebra-]{algebra}
\externaldocument[brauer-]{brauer}
\externaldocument[homology-]{homology}
\externaldocument[derived-]{derived}
\externaldocument[simplicial-]{simplicial}
\externaldocument[more-algebra-]{more-algebra}
\externaldocument[smoothing-]{smoothing}
\externaldocument[modules-]{modules}
\externaldocument[sites-modules-]{sites-modules}
\externaldocument[injectives-]{injectives}
\externaldocument[cohomology-]{cohomology}
\externaldocument[sites-cohomology-]{sites-cohomology}
\externaldocument[dga-]{dga}
\externaldocument[dpa-]{dpa}
\externaldocument[sdga-]{sdga}
\externaldocument[hypercovering-]{hypercovering}
\externaldocument[schemes-]{schemes}
\externaldocument[constructions-]{constructions}
\externaldocument[properties-]{properties}
\externaldocument[morphisms-]{morphisms}
\externaldocument[coherent-]{coherent}
\externaldocument[divisors-]{divisors}
\externaldocument[limits-]{limits}
\externaldocument[varieties-]{varieties}
\externaldocument[topologies-]{topologies}
\externaldocument[descent-]{descent}
\externaldocument[perfect-]{perfect}
\externaldocument[more-morphisms-]{more-morphisms}
\externaldocument[flat-]{flat}
\externaldocument[groupoids-]{groupoids}
\externaldocument[more-groupoids-]{more-groupoids}
\externaldocument[etale-]{etale}
\externaldocument[chow-]{chow}
\externaldocument[intersection-]{intersection}
\externaldocument[pic-]{pic}
\externaldocument[weil-]{weil}
\externaldocument[adequate-]{adequate}
\externaldocument[dualizing-]{dualizing}
\externaldocument[duality-]{duality}
\externaldocument[discriminant-]{discriminant}
\externaldocument[derham-]{derham}
\externaldocument[local-cohomology-]{local-cohomology}
\externaldocument[algebraization-]{algebraization}
\externaldocument[curves-]{curves}
\externaldocument[resolve-]{resolve}
\externaldocument[models-]{models}
\externaldocument[functors-]{functors}
\externaldocument[equiv-]{equiv}
\externaldocument[pione-]{pione}
\externaldocument[etale-cohomology-]{etale-cohomology}
\externaldocument[proetale-]{proetale}
\externaldocument[relative-cycles-]{relative-cycles}
\externaldocument[more-etale-]{more-etale}
\externaldocument[trace-]{trace}
\externaldocument[crystalline-]{crystalline}
\externaldocument[spaces-]{spaces}
\externaldocument[spaces-properties-]{spaces-properties}
\externaldocument[spaces-morphisms-]{spaces-morphisms}
\externaldocument[decent-spaces-]{decent-spaces}
\externaldocument[spaces-cohomology-]{spaces-cohomology}
\externaldocument[spaces-limits-]{spaces-limits}
\externaldocument[spaces-divisors-]{spaces-divisors}
\externaldocument[spaces-over-fields-]{spaces-over-fields}
\externaldocument[spaces-topologies-]{spaces-topologies}
\externaldocument[spaces-descent-]{spaces-descent}
\externaldocument[spaces-perfect-]{spaces-perfect}
\externaldocument[spaces-more-morphisms-]{spaces-more-morphisms}
\externaldocument[spaces-flat-]{spaces-flat}
\externaldocument[spaces-groupoids-]{spaces-groupoids}
\externaldocument[spaces-more-groupoids-]{spaces-more-groupoids}
\externaldocument[bootstrap-]{bootstrap}
\externaldocument[spaces-pushouts-]{spaces-pushouts}
\externaldocument[spaces-chow-]{spaces-chow}
\externaldocument[groupoids-quotients-]{groupoids-quotients}
\externaldocument[spaces-more-cohomology-]{spaces-more-cohomology}
\externaldocument[spaces-simplicial-]{spaces-simplicial}
\externaldocument[spaces-duality-]{spaces-duality}
\externaldocument[formal-spaces-]{formal-spaces}
\externaldocument[restricted-]{restricted}
\externaldocument[spaces-resolve-]{spaces-resolve}
\externaldocument[formal-defos-]{formal-defos}
\externaldocument[defos-]{defos}
\externaldocument[cotangent-]{cotangent}
\externaldocument[examples-defos-]{examples-defos}
\externaldocument[algebraic-]{algebraic}
\externaldocument[examples-stacks-]{examples-stacks}
\externaldocument[stacks-sheaves-]{stacks-sheaves}
\externaldocument[criteria-]{criteria}
\externaldocument[artin-]{artin}
\externaldocument[quot-]{quot}
\externaldocument[stacks-properties-]{stacks-properties}
\externaldocument[stacks-morphisms-]{stacks-morphisms}
\externaldocument[stacks-limits-]{stacks-limits}
\externaldocument[stacks-cohomology-]{stacks-cohomology}
\externaldocument[stacks-perfect-]{stacks-perfect}
\externaldocument[stacks-introduction-]{stacks-introduction}
\externaldocument[stacks-more-morphisms-]{stacks-more-morphisms}
\externaldocument[stacks-geometry-]{stacks-geometry}
\externaldocument[moduli-]{moduli}
\externaldocument[moduli-curves-]{moduli-curves}
\externaldocument[examples-]{examples}
\externaldocument[exercises-]{exercises}
\externaldocument[guide-]{guide}
\externaldocument[desirables-]{desirables}
\externaldocument[coding-]{coding}
\externaldocument[obsolete-]{obsolete}
\externaldocument[fdl-]{fdl}
\externaldocument[index-]{index}

% Theorem environments.
%
\theoremstyle{plain}
\newtheorem{theorem}[subsection]{Theorem}
\newtheorem{proposition}[subsection]{Proposition}
\newtheorem{lemma}[subsection]{Lemma}

\theoremstyle{definition}
\newtheorem{definition}[subsection]{Definition}
\newtheorem{example}[subsection]{Example}
\newtheorem{exercise}[subsection]{Exercise}
\newtheorem{situation}[subsection]{Situation}

\theoremstyle{remark}
\newtheorem{remark}[subsection]{Remark}
\newtheorem{remarks}[subsection]{Remarks}

\numberwithin{equation}{subsection}

% Macros
%
\def\lim{\mathop{\mathrm{lim}}\nolimits}
\def\colim{\mathop{\mathrm{colim}}\nolimits}
\def\Spec{\mathop{\mathrm{Spec}}}
\def\Hom{\mathop{\mathrm{Hom}}\nolimits}
\def\Ext{\mathop{\mathrm{Ext}}\nolimits}
\def\SheafHom{\mathop{\mathcal{H}\!\mathit{om}}\nolimits}
\def\SheafExt{\mathop{\mathcal{E}\!\mathit{xt}}\nolimits}
\def\Sch{\mathit{Sch}}
\def\Mor{\mathop{\mathrm{Mor}}\nolimits}
\def\Ob{\mathop{\mathrm{Ob}}\nolimits}
\def\Sh{\mathop{\mathit{Sh}}\nolimits}
\def\NL{\mathop{N\!L}\nolimits}
\def\CH{\mathop{\mathrm{CH}}\nolimits}
\def\proetale{{pro\text{-}\acute{e}tale}}
\def\etale{{\acute{e}tale}}
\def\QCoh{\mathit{QCoh}}
\def\Ker{\mathop{\mathrm{Ker}}}
\def\Im{\mathop{\mathrm{Im}}}
\def\Coker{\mathop{\mathrm{Coker}}}
\def\Coim{\mathop{\mathrm{Coim}}}

% Boxtimes
%
\DeclareMathSymbol{\boxtimes}{\mathbin}{AMSa}{"02}

%
% Macros for moduli stacks/spaces
%
\def\QCohstack{\mathcal{QC}\!\mathit{oh}}
\def\Cohstack{\mathcal{C}\!\mathit{oh}}
\def\Spacesstack{\mathcal{S}\!\mathit{paces}}
\def\Quotfunctor{\mathrm{Quot}}
\def\Hilbfunctor{\mathrm{Hilb}}
\def\Curvesstack{\mathcal{C}\!\mathit{urves}}
\def\Polarizedstack{\mathcal{P}\!\mathit{olarized}}
\def\Complexesstack{\mathcal{C}\!\mathit{omplexes}}
% \Pic is the operator that assigns to X its picard group, usage \Pic(X)
% \Picardstack_{X/B} denotes the Picard stack of X over B
% \Picardfunctor_{X/B} denotes the Picard functor of X over B
\def\Pic{\mathop{\mathrm{Pic}}\nolimits}
\def\Picardstack{\mathcal{P}\!\mathit{ic}}
\def\Picardfunctor{\mathrm{Pic}}
\def\Deformationcategory{\mathcal{D}\!\mathit{ef}}
\usepackage{cleveref}
\crefname{originallemma}{original lemma}{original lemmas}
\Crefname{originallemma}{Original lemma}{Original lemmas}
\makeatletter
\newlabel{algebra-lemma-cover-module}{{00EK}{source}{}{sourcepointer-1}{}}
\newlabel{algebra-lemma-cover-module@cref}{{[originallemma][0][]00EK}{[1][0][]source}{}{}{}}
\makeatother
\crefname{originalequation}{original equation}{original equations}
\Crefname{originalequation}{Original equation}{Original equations}
\makeatletter
\newlabel{equation-alternative-ringed}{{0GZU}{source}{}{sourcepointer-2}{}}
\newlabel{equation-alternative-ringed@cref}{{[originalequation][0][]0GZU}{[1][0][]source}{}{}{}}
\makeatother
\crefname{originalproposition}{original proposition}{original propositions}
\Crefname{originalproposition}{Original proposition}{Original propositions}
\makeatletter
\newlabel{proposition-equivalence-quasi-coherent}{{03DX}{source}{}{sourcepointer-3}{}}
\newlabel{proposition-equivalence-quasi-coherent@cref}{{[originalproposition][0][]03DX}{[1][0][]source}{}{}{}}
\makeatother
\crefname{originallemma}{original lemma}{original lemmas}
\Crefname{originallemma}{Original lemma}{Original lemmas}
\makeatletter
\newlabel{schemes-lemma-widetilde-pullback}{{01I9}{source}{}{sourcepointer-4}{}}
\newlabel{schemes-lemma-widetilde-pullback@cref}{{[originallemma][0][]01I9}{[1][0][]source}{}{}{}}
\makeatother
\crefname{originalsection}{original section}{original sections}
\Crefname{originalsection}{Original section}{Original sections}
\makeatletter
\newlabel{section-quasi-coherent-sheaves}{{03DR}{source}{}{sourcepointer-5}{}}
\newlabel{section-quasi-coherent-sheaves@cref}{{[originalsection][0][]03DR}{[1][0][]source}{}{}{}}
\makeatother
\crefname{originaldefinition}{original definition}{original definitions}
\Crefname{originaldefinition}{Original definition}{Original definitions}
\makeatletter
\newlabel{sites-modules-definition-site-local}{{03DL}{source}{}{sourcepointer-6}{}}
\newlabel{sites-modules-definition-site-local@cref}{{[originaldefinition][0][]03DL}{[1][0][]source}{}{}{}}
\makeatother
\crefname{originallemma}{original lemma}{original lemmas}
\Crefname{originallemma}{Original lemma}{Original lemmas}
\makeatletter
\newlabel{sites-modules-lemma-special-locally-free}{{03DM}{source}{}{sourcepointer-7}{}}
\newlabel{sites-modules-lemma-special-locally-free@cref}{{[originallemma][0][]03DM}{[1][0][]source}{}{}{}}
\makeatother
\title{Eight equivalent descriptions of quasi-coherent modules on the category of affine schemes}
\author{The Stacks Project Authors}
\date{stacks-project:descent.tex:local-snapshot-20261012; GNU FDL 1.2 or later}
\begin{document}
\maketitle
\noindent\textit{Editorial note.} All delivered mathematical text is the author's own, copied from the pinned Stacks Project chapter file descent.tex; no statement or proof step was rewritten, reordered or removed.\par\medskip
\noindent\textit{Editorial note.} The author's own closing aside 'With apologies for the horrible notation' is quoted verbatim as part of the proof.\par\medskip
\noindent\textit{Editorial note.} The counted spans carry navigation stubs for the references leaving them; no bibliography citation occurs in the spans.\par\medskip
\noindent\textit{Editorial note.} The preamble span points at the chapter's own preamble.tex shipped with the frozen source, used verbatim as staged.\par\medskip
\noindent\textit{Allowed standard prerequisites (editorial).}\begin{itemize}
\item Quasi-coherence is intrinsic, so the equivalence of ringed sites preserves quasi-coherent modules. Modules on Sites, Lemma sites-modules-lemma-special-locally-free, invoked at proof lines 3779-3782. Transfers quasi-coherence along the equivalence (equation-alternative-ringed).
\item Topology-independence of quasi-coherent modules on Sch/S and identification with the O\_S-module procedure. Same chapter, Proposition proposition-equivalence-quasi-coherent, invoked at proof lines 3783-3786. Shows (2)-(7) are equivalent and that (8) is the same as (7).
\item The sheaf condition for a module determined by localisations on a standard cover. Algebra, Lemma algebra-lemma-cover-module, invoked at proof lines 3814-3815. Gives that F is a Zariski sheaf once the base-change formula is available.
\item Compatibility of pullback with the tilde construction, and right exactness of the tensor product. Schemes, Lemma schemes-lemma-widetilde-pullback, invoked at proof lines 3795-3795, and the right exactness of tensor product used at proof lines 3825-3835. Gives (1) implies (2) and the local presentations of F.
\end{itemize}\medskip
\ifdefined\hypertarget\hypertarget{sourcepointer-1}{}\fi
\noindent\textit{Source reference.} Original lemma 00EK: sources/stacks-chow/tags.txt, source lines 537--537. Pointer outside the counted span, invoked once for context; the load-bearing argument of the counted statement is the author's proof reproduced in full in the proof span.\par\medskip
\ifdefined\hypertarget\hypertarget{sourcepointer-2}{}\fi
\noindent\textit{Source reference.} Original equation 0GZU: sources/stacks-algebraic-stacks-sheaves/pinned-original-source.tex, source lines 4727--4727. Pointer outside the counted span, invoked once for context; the load-bearing argument of the counted statement is the author's proof reproduced in full in the proof span.\par\medskip
\ifdefined\hypertarget\hypertarget{sourcepointer-3}{}\fi
\noindent\textit{Source reference.} Original proposition 03DX: sources/stacks-descent/pinned-original-source.tex, source lines 2716--2716. Pointer outside the counted span, invoked once for context; the load-bearing argument of the counted statement is the author's proof reproduced in full in the proof span.\par\medskip
\ifdefined\hypertarget\hypertarget{sourcepointer-4}{}\fi
\noindent\textit{Source reference.} Original lemma 01I9: sources/stacks-chow/tags.txt, source lines 1970--1970. Pointer outside the counted span, invoked once for context; the load-bearing argument of the counted statement is the author's proof reproduced in full in the proof span.\par\medskip
\ifdefined\hypertarget\hypertarget{sourcepointer-5}{}\fi
\noindent\textit{Source reference.} Original section 03DR: sources/stacks-descent/pinned-original-source.tex, source lines 2355--2355. Pointer outside the counted span, invoked once for context; the load-bearing argument of the counted statement is the author's proof reproduced in full in the proof span.\par\medskip
\ifdefined\hypertarget\hypertarget{sourcepointer-6}{}\fi
\noindent\textit{Source reference.} Original definition 03DL: sources/stacks-chow/tags.txt, source lines 4424--4424. Pointer outside the counted span, invoked once for context; the load-bearing argument of the counted statement is the author's proof reproduced in full in the proof span.\par\medskip
\ifdefined\hypertarget\hypertarget{sourcepointer-7}{}\fi
\noindent\textit{Source reference.} Original lemma 03DM: sources/stacks-chow/tags.txt, source lines 4425--4425. Pointer outside the counted span, invoked once for context; the load-bearing argument of the counted statement is the author's proof reproduced in full in the proof span.\par\medskip
\begin{lemma}
\label{lemma-quasi-coherent-alternative}
Let $S$ be a scheme. Let $\mathcal{F}$ be a presheaf
of $\mathcal{O}_{\textit{Aff}}$-modules on $(\textit{Aff}/S)_{fppf}$.
The following are equivalent
\begin{enumerate}
\item for every morphism $U \to U'$ of $(\textit{Aff}/S)_{fppf}$ the map
$\mathcal{F}(U') \otimes_{\mathcal{O}(U')} \mathcal{O}(U) \to \mathcal{F}(U)$
is an isomorphism,
\item $\mathcal{F}$ is a sheaf on $(\textit{Aff}/S)_{Zar}$ and
a quasi-coherent module on the ringed site
$((\textit{Aff}/S)_{Zar}, \mathcal{O}_{\textit{Aff}})$ in the sense of
Modules on Sites, Definition \ref{sites-modules-definition-site-local},
\item same as in (2) for the \'etale topology,
\item same as in (2) for the smooth topology,
\item same as in (2) for the syntomic topology,
\item same as in (2) for the fppf topology,
\item $\mathcal{F}$ corresponds to a quasi-coherent module on
$(\Sch/S)_{Zar}$,
$(\Sch/S)_\etale$,
$(\Sch/S)_{smooth}$,
$(\Sch/S)_{syntomic}$, or
$(\Sch/S)_{fppf}$
via the equivalence (\ref{equation-alternative-ringed}),
\item $\mathcal{F}$ comes from a unique quasi-coherent
$\mathcal{O}_S$-module $\mathcal{G}$ by the procedure
described in Section \ref{section-quasi-coherent-sheaves}.
\end{enumerate}
\end{lemma}

\begin{proof}
Since the notion of a quasi-coherent module is intrinsic
(Modules on Sites, Lemma \ref{sites-modules-lemma-special-locally-free})
we see that the equivalence (\ref{equation-alternative-ringed})
induces an equivalence between categories of quasi-coherent modules.
Proposition \ref{proposition-equivalence-quasi-coherent}
says the topology we use to study quasi-coherent modules on
$\Sch/S$ does not matter and it also tells us that (8)
is the same as (7). Hence we see that (2) -- (8) are all equivalent.

\medskip\noindent
Assume the equivalent conditions (2) -- (8) hold and let
$\mathcal{G}$ be as in (8). Let $h : U \to U' \to S$ be a morphism
of $\textit{Aff}/S$. Denote $f : U \to S$ and $f' : U' \to S$ the
structure morphisms, so that $f = f' \circ h$.
We have $\mathcal{F}(U') = \Gamma(U', (f')^*\mathcal{G})$ and
$\mathcal{F}(U) = \Gamma(U, f^*\mathcal{G}) = \Gamma(U, h^*(f')^*\mathcal{G})$.
Hence (1) holds by Schemes, Lemma \ref{schemes-lemma-widetilde-pullback}.

\medskip\noindent
Assume (1) holds. To finish the proof it suffices to prove (2).
Let $U$ be an object of $(\textit{Aff}/S)_{Zar}$.
Say $U = \Spec(R)$. A standard open covering $U = U_1 \cup \ldots \cup U_n$
is given by $U_i = D(f_i)$ for some elements $f_1, \ldots, f_n \in R$
generating the unit ideal of $R$. By property (1) we see that
$$
\mathcal{F}(U_i) =
\mathcal{F}(U) \otimes_R R_{f_i} =
\mathcal{F}(U)_{f_i}
$$
and
$$
\mathcal{F}(U_i \cap U_j) =
\mathcal{F}(U) \otimes_R R_{f_if_j} =
\mathcal{F}(U)_{f_if_j}
$$
Thus we conclude from Algebra, Lemma \ref{algebra-lemma-cover-module}
that $\mathcal{F}$ is a sheaf on $(\textit{Aff}/S)_{Zar}$. Choose a
presentation
$$
\bigoplus\nolimits_{k \in K} R
\longrightarrow
\bigoplus\nolimits_{l \in L} R
\longrightarrow
\mathcal{F}(U)
\longrightarrow 0
$$
by free $R$-modules. By property (1) and the right exactness of tensor product
we see that for every morphism $U' \to U$ in $(\textit{Aff}/S)_{Zar}$
we obtain a presentation
$$
\bigoplus\nolimits_{k \in K} \mathcal{O}_{Aff}(U')
\longrightarrow
\bigoplus\nolimits_{l \in L} \mathcal{O}_{Aff}(U')
\longrightarrow
\mathcal{F}(U')
\longrightarrow 0
$$
In other words, we see that the restriction of $\mathcal{F}$
to the localized category $(\textit{Aff}/S)_{Zar}/U$ has a presentation
$$
\bigoplus\nolimits_{k \in K} \mathcal{O}_{Aff}|_{(\textit{Aff}/S)_{Zar}/U}
\longrightarrow
\bigoplus\nolimits_{l \in L} \mathcal{O}_{Aff}|_{(\textit{Aff}/S)_{Zar}/U}
\longrightarrow
\mathcal{F}|_{(\textit{Aff}/S)_{Zar}/U}
\longrightarrow 0
$$
With apologies for the horrible notation, this finishes the proof.
\end{proof}

\end{document}
